% Laget automatisk av elevark/lag_elevark.py fra opplegg/repetisjon-krefter.html. Endringer her blir overskrevet.
\documentclass[11pt]{article}
\usepackage{elevark}
\begin{document}
\elevarkhode{Fysikk 1 · Krefter}{Repetisjon: krefter}{https://lektorbo.no/opplegg/repetisjon-krefter.html}{Opplegget på nett}

Med Newtons lover kan vi regne ut bevegelsen når vi kjenner kreftene, og omvendt. Her er en oversikt over lovene og kreftene, og blandede oppgaver med løsningsforslag.

\section*{Newtons lover}

\begin{center}\small
\begin{tabularx}{\linewidth}{YY}
\toprule
Første lov & Når summen av kreftene er null, er gjenstanden i ro eller har konstant fart i en rett linje. \\
Andre lov & $ΣF = ma$. Kraftsummen og akselerasjonen har samme retning. \\
Tredje lov & Når $A$ virker på $B$ med en kraft, virker $B$ på $A$ med en like stor kraft i motsatt retning. Kraft og motkraft virker på forskjellige gjenstander og opphever ikke hverandre. \\
\bottomrule
\end{tabularx}
\end{center}

\section*{Slik løser du en kraftoppgave}

\begin{enumerate}
\item Velg én gjenstand, og tegn bare kreftene som virker på den.
\item Velg en positiv retning, gjerne i akselerasjonsretningen. På et skråplan: langs og vinkelrett på planet.
\item Skriv Newtons andre lov for hver retning, og løs likningene.
\end{enumerate}

\section*{Formler}

\begin{itemize}
\item Tyngde: $G = mg$. Friksjon ved glidning: $R = μN$. Hvilefriksjonen er høyst $μ_{\mathrm{s}}N$.
\item Luftmotstand: $L = kv$ (lav fart) eller $L = kv^{2}$ (høy fart). Terminalfart når $L = mg$.
\item Skråplan med vinkel $θ$: tyngdekomponenten langs planet er $mg \sin  θ$, og vinkelrett på planet $mg \cos  θ$. Uten andre krefter vinkelrett på planet er $N = mg \cos  θ$.
\item Heis: $N = m(g + a)$, med $a$ positiv oppover.
\end{itemize}

\section*{Oppgaver}

Bruk $g = 9{,}81\,\mathrm{m/s^{2}}$.

\begin{oppgaveliste}
\oppgave \fase{forklar}{Begrep}Er påstandene sanne eller usanne? Begrunn svaret.

a) Normalkraften er alltid like stor som tyngden.\newline b) Kraft og motkraft opphever hverandre.\newline c) En bil som kjører med konstant fart i en rett linje, har kraftsum null.\newline d) Friksjonskraften peker alltid motsatt av fartsretningen.\newline e) En tung og en lett stein faller like fort uten luftmotstand.
\svar
\oppgave \fase{test}{Regn}En person på 70~kg står på en badevekt i en heis som akselererer oppover med 1,2~m/s². Finn normalkraften og hva vekta viser. Sjekk med simuleringen.\footnote{\url{lektorbo.no/sim/heis.html}}
\svar
\oppgave \fase{test}{Regn}En kasse på 6,0~kg trekkes bortover et gulv med en vannrett kraft på 40~N. Friksjonstallet er 0,30. Finn akselerasjonen. Sjekk med simuleringen.\footnote{\url{lektorbo.no/sim/friksjon.html}}
\svar
\oppgave \fase{test}{Regn}En kloss på 2,0~kg glir ned et skråplan med helningsvinkel 30°. Friksjonstallet er 0,20. Finn normalkraften, friksjonskraften og akselerasjonen. Sjekk med simuleringen.\footnote{\url{lektorbo.no/sim/skraaplan.html}}
\svar
\oppgave \fase{test}{Regn}En vogn på 0,80~kg står på et bord uten friksjon. Den er koblet med en snor over en trinse til et lodd på 0,20~kg som henger fritt. Finn akselerasjonen og snorkraften. Sjekk med simuleringen.\footnote{\url{lektorbo.no/sim/vogn-og-lodd.html}}
\svar
\oppgave \fase{test}{Regn}En fallskjermhopper på 80~kg har terminalfarten 55~m/s før skjermen åpnes og 5,0~m/s etterpå. Luftmotstanden er $L = kv^{2}$. Finn $k$ før og etter at skjermen åpnes. Hvor stor er akselerasjonen når farten er 30~m/s, før skjermen åpnes? Se Fall med luftmotstand.\footnote{\url{lektorbo.no/sim/luftmotstand.html}}
\svar
\oppgave \fase{utfordring}{Sammensatt}En bil på 1200~kg i 90~km/h bråbremser og stopper på 40~m. Finn bremsekraften. Hvor stort må friksjonstallet mellom dekk og vei minst være? Hva skjer med bremselengden på glatt føre med $μ = 0{,}20$?
\svar
\end{oppgaveliste}

\end{document}
